% =============================================================================
% Section 1: Introduction & Contribution
% =============================================================================
\section{Introduction and Contribution}\label{sec:introduction}

\subsection{Motivation}

Composite indicators that measure structural features of national
economies have a long tradition in institutional research.  Indices
published by the OECD, the World Economic Forum, and various United
Nations agencies compress multidimensional data into scalar summaries
that facilitate cross-country comparison and policy
monitoring~\cite{oecd_handbook_2008,saisana2005}.  The International
Sovereignty Index (ISI) contributes to this tradition by focusing
specifically on \emph{concentration of bilateral external dependencies}:
the degree to which a country's cross-border economic and strategic
linkages are concentrated among a small number of counterparty states.

The intellectual antecedents of the ISI lie in two distinct literatures.
First, the industrial-organisation literature on market concentration,
where the Herfindahl--Hirschman Index (HHI) quantifies the degree to
which market shares are concentrated among a small number of
firms~\cite{herfindahl1950,hirschman1964}.  Second, the literature on
supply-chain vulnerability, which has expanded substantially since the
semiconductor shortages of 2020--2021 and the geopolitical
disruptions following the 2022 Russia--Ukraine
conflict~\cite{ec_critical_raw}.  The ISI bridges these traditions by
applying concentration measurement to bilateral trade and financial
flows across six strategically distinct domains.

The ISI does not measure sovereignty in a political or legal sense.
Rather, it quantifies a structural feature---counterparty
concentration---that is a prerequisite for informed policy analysis of
external dependence.  A high ISI score indicates that a country's
external flows in a given domain are heavily concentrated among few
partners; a low score indicates dispersion across many partners.  The
index makes no normative claim about whether concentration is
advantageous or disadvantageous.  The question of whether high
concentration translates into political vulnerability, economic
fragility, or coercive leverage is a distinct empirical question that
requires analysis beyond the ISI's measurement scope.

\subsection{Scope of this paper}

This paper---the second in the ISI series---presents the full
empirical results of the first ISI application: 27~EU member states,
vintage-2024 data, methodology v1.0.  The methodological foundations
(axis definitions, normalisation procedures, data sources, aggregation
rule, integrity chain) are documented in Paper~1~\cite{isi_paper1}.
The present paper is self-contained for readers interested in the
empirical findings; cross-references to Paper~1 are provided where
construction details are relevant to interpretation.

Relative to a conventional ``results'' paper in the composite-indicator
literature, this paper provides substantially deeper analytical content.
Beyond reporting the ranking and basic distributional statistics, we
undertake a formal variance decomposition with proof, systematic
robustness exercises with both Spearman and Kendall rank-order
statistics, rank-elasticity analysis, axis contribution decomposition,
classification compression diagnostics, structural analysis of the
defence axis, compensability analysis, and geographic correlate
tabulation.  Each analytical component is designed to interrogate the
ISI composite from a different angle and to provide the reader with
the evidence needed to form an independent judgment about the index's
informational content.

The specific contributions of this paper are:

\begin{enumerate}[leftmargin=2em]
  \item A complete composite and axis-level ranking of all 27~EU
    member states with eight-decimal precision, accompanied by
    classification labels and distributional diagnostics (histogram,
    ECDF, Lorenz curve, Gini coefficient).
  \item A formal variance decomposition---proved analytically from the
    linear aggregation rule---decomposing the composite variance into
    own-variance and cross-covariance components, with percentage
    contributions for each axis.
  \item A leave-one-axis-out (LOO) sensitivity analysis with both
    Spearman~$\rho$ and Kendall~$\tau$, accompanied by a full $27
    \times 6$ rank-displacement matrix.
  \item A rank-elasticity analysis quantifying the sensitivity of each
    country's rank to marginal perturbations in individual axis scores.
  \item An axis-contribution analysis providing bivariate $R^2$ and
    partial correlations between each axis and the composite rank.
  \item Classification compression diagnostics comparing the
    official four-tier scheme with data-driven tercile and quartile
    alternatives, reporting within-tier variance and $\eta^2$.
  \item A structural analysis of the defence axis examining its
    distributional properties, oligopsonistic character, and
    country-level defence-to-max-axis gaps.
  \item A cross-axis correlation analysis with formal $t$-statistics,
    eigenvalue decomposition, and network representation.
  \item A compensability analysis documenting the degree to which the
    arithmetic-mean aggregation rule permits cross-axis compensation.
  \item An assessment of geographic and structural correlates (island
    status, population size, core/periphery classification).
  \item A four-part robustness programme: LOO, z-score
    standardisation, Dirichlet weight perturbation (Monte Carlo), and
    winsorisation.
  \item Radar-chart profiles for eight selected countries illustrating
    the heterogeneity of axis-level patterns.
  \item A reproducibility protocol and computational appendix enabling
    full replication from the provided snapshot.
\end{enumerate}

\subsection{Interpretation conventions}

Throughout this paper, the following conventions apply:

\begin{itemize}[leftmargin=2em]
  \item Higher ISI composite $\Rightarrow$ higher concentration
    $\Rightarrow$ greater structural dependence on a small set of
    counterparties.
  \item Rank~1 = most concentrated country; rank~27 = least
    concentrated.
  \item The terms ``most concentrated'' and ``least concentrated'' are
    descriptive.  They are not synonymous with ``worst'' and ``best.''
  \item Correlation between axes does not imply causation.  Where
    plausible structural interpretations exist, they are stated as
    hypotheses, not conclusions.
  \item Population standard deviation ($N$ denominator) is used
    throughout because the 27~EU member states constitute the full
    target population, not a sample drawn from a larger universe.
  \item Statistical significance of correlation coefficients is
    assessed via the $t$-test: $t = r\sqrt{(N-2)/(1-r^2)}$, with
    $N = 27$, df~$= 25$, and a two-sided critical value of
    $t_{0.025,25} = 2.060$ at the 5\% level.
\end{itemize}

\subsection{Structure}

The remainder of this paper is organised as follows.
\Cref{sec:data-construction} summarises the data sources, aggregation
rule, and classification thresholds.
\Cref{sec:composite-results} presents the composite distribution,
full ranking, tier-compression analysis, and classification
compression diagnostics.
\Cref{sec:axis-diagnostics} provides axis-level diagnostics including
dispersion, dominance, variance decomposition with formal proof,
rank contribution analysis, and defence axis structural analysis.
\Cref{sec:cross-axis} examines the pairwise correlation structure
with $t$-statistics and eigenvalue decomposition.
\Cref{sec:robustness} reports the robustness programme: LOO with
Kendall $\tau$, rank-displacement matrices, z-score standardisation,
Dirichlet weight perturbation, winsorisation, and rank-elasticity
analysis.
\Cref{sec:compensability} presents the compensability analysis.
\Cref{sec:country-profiles} presents radar-chart profiles for eight
selected countries and geographic correlates.
\Cref{sec:discussion} discusses interpretation boundaries and
limitations.
\Cref{sec:conclusion} concludes.
Appendices~A--D provide full tables, robustness details, the
reproducibility protocol, and the computational code listing.
